%% kapa: exam=1 round=37 year=2026 period=1 form=A subject=civil_law subject_ko=민법 seq=4 number=4
%% source: Q-Net 감정평가사 기출문제 (cst00309) artlSeq=5258924 file=raw/1008/5258924/2026년 제37회 감정평가사 1차 1교시.pdf
%% conversion: kordoc 4.24.51 -> tools/gen1.py (auto); review=proofread-sonnet; flags=none
\begin{question}{4}
甲이 생전처분으로 광주에 소재하는 그 소유 X부동산을 출연하여 서울에 주된 사무소를 두는 민법상 재단법인(A)을 설립하고자 한다. 이에 관한 설명으로 옳은 것을 모두 고른 것은? (다툼이 있으면 판례에 따름)
\begin{framedtext}
ㄱ. A는 X부동산의 소재지인 광주에서 설립등기를 함으로써 성립한다. \par
ㄴ. 甲은 사원자격의 득실에 관한 사항을 정관에 기재하고 이를 등기하여야 한다. \par
ㄷ. 甲이 이사임면의 방법을 정하지 않고 사망한 경우, 이해관계인 또는 검사의 청구에 의하여 법원이 이를 정한다. \par
ㄹ. 甲은 A가 성립한 이후에는 착오에 의한 의사표시의 취소 요건을 갖추더라도 이를 이유로 그 출연행위를 취소할 수 없다. \par
\end{framedtext}
\begin{choices}
  \item ㄷ
  \item ㄱ, ㄴ
  \item ㄷ, ㄹ
  \item ㄱ, ㄴ, ㄹ
  \item ㄱ, ㄴ, ㄷ, ㄹ
\end{choices}
\end{question}
\answer{1}
